% Pista alg-23s-q1 · Àlgebra
% Procedència: PAU Catalunya 2023 · Convocatòria de setembre · Sèrie 2 · Qüestió 1
\documentclass[11pt,a4paper]{article}
\usepackage{pista}
\pistainfo{Catalunya 2023}{Convocatòria de setembre}{Sèrie 2}

\begin{document}

\section*{\textcolor{blauPista}{Qüestió 1}}

\noindent Siguin
\[
A = \begin{pmatrix} 2 & 1 \\ 3 & 2 \end{pmatrix}, \quad
B = \begin{pmatrix} 2 & -1 \\ -3 & 2 \end{pmatrix} \quad \text{i la matriu identitat d'ordre dos } I = \begin{pmatrix} 1 & 0 \\ 0 & 1 \end{pmatrix}.
\]

\subsection*{\textit{a)} \normalfont Comproveu que $(A - 2I)^{2} = 3I$. \hfill\textnormal{[0,5~punts]}}

\pista{És un càlcul directe. Comença obtenint la matriu $A - 2I$ (resta element a element), eleva-la al quadrat, i comprova que el resultat coincideix amb $3I$. Recorda que elevar al quadrat una matriu significa multiplicar-la per ella mateixa, i la multiplicació de matrius ja l'has estudiat.}

\subsection*{\textit{b)} \normalfont Utilitzant la igualtat de l'apartat anterior, trobeu la matriu inversa de la matriu $A$ en funció de les matrius $A$ i $I$, i comproveu que coincideix amb la matriu $B$. \hfill\textnormal{[1,25~punts]}}

\pista{La idea clau d'aquest apartat és la següent: si aconsegueixes manipular la igualtat $(A - 2I)^{2} = 3I$ fins a portar-la a la forma $A \cdot M = I$, on $M$ és una expressió en funció d'$A$ i $I$, llavors aquesta $M$ serà, per definició, la matriu inversa $A^{-1}$. Per fer-ho, comença desenvolupant el quadrat $(A - 2I)^{2}$ i recorda que $I$ es comporta com l'1 en el producte de matrius (per exemple, $A \cdot I = A$). Apareixerà $A^{2}$, que pots reescriure com $A \cdot A$ per poder ``treure factor comú $A$'' a la part de l'esquerra. Quan tinguis $M$, calcula-la numèricament i comprova que dóna exactament $B$.}

\subsection*{\textit{c)} \normalfont Calculeu la matriu $X$ que satisfà la igualtat $A \cdot X = B$. \hfill\textnormal{[0,75~punts]}}

\pista{A l'apartat b) ja has descobert que $A^{-1} = B$, així que ja tens la inversa d'$A$ feta i no l'has de tornar a calcular. Per aïllar $X$ de la igualtat $A \cdot X = B$, multiplica per $A^{-1}$ a banda i banda, per l'esquerra. Quan substitueixis, veuràs que $X$ es redueix al producte de dues matrius que ja coneixes.}

\end{document}
